% !TEX program = xelatex
\documentclass[11pt, a4paper]{cvClass}

% Configure page margins
\geometry{left=2cm, top=1.2cm, right=2cm, bottom=1.2cm, footskip=.5cm}

% Color theme
\colorlet{awesome}{awesome-red}

% Uncomment if you want a photo (and add file)
\photo[circle,edge,left]{Tobias_Staal_square}
\name{Tobias}{Stål}
\position{Geoscientist and engineer}
\address{Tasmania, Australia}

% Contact
\mobile{+61 4 4749 4089}
\email{tobias.staal@utas.edu.au}
\homepage{tobbetripitaka.github.io/Tobias\_Staal}
\github{TobbeTripitaka}

\begin{document}

% Header
\makecvheader

\hspace{2cm}

%===============================================================================
% Peer-reviewed publications
%===============================================================================
\cvsection{Peer-reviewed publications}

\cvsubsection{Overview}
\begin{cvparagraph}
Peer-reviewed journal articles in chronological order, including accepted papers and manuscripts currently under review.
\end{cvparagraph}

\cvsubsection{Selected publications}
\begin{cvhonors}
  \cvhonor
    {2026}
    {Halpin, J. A., N. R. Daczko, J. A. Mulder, A. Maritati, T. Stål — Ediacaran–Cambrian high-strain melt-present deformation of the Archean Charcot Province, East Antarctica. \textit{Tectonics}. doi:10.1029/2026TC009479}
    {}
    {}
  \cvhonor
    {2026}
    {Manassero, M. C., K. Selway, A. M. Reading, N. Askey-Doran, J. R. Peacock, F. Ramirez, T. Stål — Magnetotellurics, mantle viscosity, integrated geophysics, glacial isostatic adjustment, mantle water content, East Antarctica. \textit{Journal of Geophysical Research: Solid Earth}. In review.}
    {}
    {}
  \cvhonor
    {2026}
    {Manassero, M. C., K. Selway, T. Stål, M. Scheiter, F. S. McCormack, M. Lösing, J. A. Halpin, B. Kulessa, A. M. Reading — Bed topography and subglacial conditions of Denman Glacier, East Antarctica: insights from magnetotelluric data and interdisciplinary studies. \textit{Journal of Geophysical Research: Earth Surface}. doi:10.1029/2026JF009277}
    {}
    {}
  \cvhonor
    {2026}
    {Kelly, I. D., A. M. Reading, T. Stål, B. Kulessa, A. García-Jerez, J. Piña-Flores, E. Paolucci, A. Tanzini, R. J. Turner, J. C. Magyar, A. P. Bassom — Determining the character of subglacial sediments in the ice–bedrock interface zone of Antarctica using horizontal-to-vertical spectral ratios (HVSRs) of seismic ambient noise. \textit{Journal of Geophysical Research: Solid Earth}. doi:10.1029/2025JB033029}
    {}
    {}
  \cvhonor
    {2026}
    {McCormack, F. S., T. Stål, N. Shao, E. MacKie, A. Fabela Hinojosa, M. Lösing, J. Roberts, S. Ehrenfeucht, C. Dow — Synthetic bed topographies for Antarctica and their utility in ice-sheet modelling. \textit{Philosophical Transactions of the Royal Society A}. doi:10.1098/rsta.2024.0537}
    {}
    {}
  \cvhonor
    {2026}
    {Arora, D., N. C. Pant, T. Stål, A. Naik, M. Pandey, R. Gupta — Geothermal heat flow in the Princess Elizabeth Land sector, East Antarctica. \textit{Earth-Science Reviews}. doi:10.1016/j.earscirev.2026.105475}
    {}
    {}
  \cvhonor
    {2026}
    {Al-Aghbary, M., M. Sobh, T. Stål, M. O. Awaleh, M. Jalludin, C. Gerhards — Improving geothermal heat-flow predictions and uncertainty quantification using clustering-based quantile regression forests. \textit{Geophysical Journal International}. doi:10.1093/gji/ggag1131}
    {}
    {}
  \cvhonor
    {2026}
    {Lösing, M., W. Colgan, T. Stål, J. Ebbing, A. G. Busck, T. Zhang, H. Seroussi, F. McCormack, D. Fahrner, et al. — Community heat-flow recommendations: suitable basal boundary conditions for Greenland and Antarctica in ISMIP7. \textit{GEUS Bulletin}. doi:10.34194/r0w9rf81}
    {}
    {}
  \cvhonor
    {2025}
    {Lösing, M., A. R. A. Aitken, J. Ebbing, J. A. Halpin, L. Li, M. Moorkamp, A. M. Reading, T. Stål — Linking tectonics and crustal thermal properties in southwestern Australia and East Antarctica through coupled gravity and magnetic analysis. \textit{Journal of Geophysical Research: Solid Earth}. doi:10.1029/2024JB030770}
    {}
    {}
  \cvhonor
    {2025}
    {Abram, N. J., A. Purich, M. H. England, F. S. McCormack, J. M. Strugnell, D. M. Bergstrom, T. R. Vance, T. Stål, B. Wienecke, P. Heil et al. — Emerging evidence of abrupt changes in the Antarctic environment. \textit{Nature}, 644(8077), pp. 621–633. doi:10.1038/s41586-025-09349-5}
    {}
    {}
    \end{cvhonors}
\begin{cvhonors}
  \cvhonor
    {2024}
    {Stål, T., J. A. Halpin, J. W. Goodge, A. M. Reading — Geology matters for Antarctic geothermal heat. \textit{Geophysical Research Letters}, 51(13), e2024GL110098. doi:10.1029/2024GL110098}
    {}
    {}
  \cvhonor
    {2023}
    {Fuchs, S., F. Neumann, B. Norden, G. Beardsmore, P. Chiozzi, W. Colgan, A. P. Dominguez, \dots, T. Stål et al. — The Global Heat Flow Database: Update 2023. \textit{GFZ Data Services}. doi:10.5880/fidgeo.2023.023}
    {}
    {}

  \cvhonor
    {2022}
    {McCormack, F. S., J. L. Roberts, C. F. Dow, T. Stål, J. A. Halpin, A. M. Reading, M. J. Siegert — Fine-scale geothermal heat flow in Antarctica can increase simulated subglacial melt estimates. \textit{Geophysical Research Letters}, 49(15), e2022GL098539. doi:10.1029/2022GL098539}
    {}
    {}
  \cvhonor
    {2022}
    {Stål, T., A. M. Reading, S. Fuchs, J. A. Halpin, M. Lösing, R. J. Turner — Properties and biases of the global heat-flow compilation. \textit{Frontiers in Earth Science}, 10, 963525. doi:10.3389/feart.2022.963525}
    {}
    {}
  \cvhonor
    {2022}
    {Reading, A. M., T. Stål, J. A. Halpin, M. Lösing, J. Ebbing, W. Shen, F. S. McCormack, C. S. Siddoway, D. Hasterok — Antarctic geothermal heat flow and its implications for tectonics and ice sheets. \textit{Nature Reviews Earth \& Environment}, 3(12), pp. 814–831. doi:10.1038/s43017-022-00348-y}
    {}
    {}
  \cvhonor
    {2021}
    {Stål, T., A. M. Reading, J. A. Halpin, J. M. Whittaker — Antarctic geothermal heat-flow model: Aq1. \textit{Geochemistry, Geophysics, Geosystems}, 22(2), e2020GC009428. doi:10.1029/2020GC009428}
    {}
    {}
  \cvhonor
    {2021}
    {Sanchez, G., J. A. Halpin, M. Gard, D. Hasterok, T. Stål, T. Raimondo, S. Peters, A. Burton-Johnson — PetroChron Antarctica: a geological database for interdisciplinary use. \textit{Geochemistry, Geophysics, Geosystems}, 22(12), e2021GC010154. doi:10.1029/2021GC010154}
    {}
    {}
  \cvhonor
    {2020}
    {Stål, T., A. M. Reading, J. A. Halpin, S. J. Phipps, J. M. Whittaker — The Antarctic crust and upper mantle: a flexible 3D model and software framework for interdisciplinary research. \textit{Frontiers in Earth Science}, 8, 577502. doi:10.3389/feart.2020.577502}
    {}
    {}
  \cvhonor
    {2020}
    {Morse, P. E., A. M. Reading, T. Stål — Exploratory volumetric deep Earth visualisation by 2.5D interactive compositing. \textit{IEEE Transactions on Visualization and Computer Graphics}, 28(7), pp. 2641–2653. doi:10.1109/TVCG.2020.3037226}
    {}
    {}
  \cvhonor
    {2020}
    {Stål, T., A. M. Reading — A grid for multidimensional and multivariate spatial representation and data processing. \textit{Journal of Open Research Software}, 8(1), pp. 1–7. doi:10.5334/jors.287}
    {}
    {}
  \cvhonor
    {2019}
    {Stål, T., A. M. Reading, J. A. Halpin, J. M. Whittaker — A multivariate approach for mapping lithospheric domain boundaries in East Antarctica. \textit{Geophysical Research Letters}, 46(17–18), pp. 10404–10416. doi:10.1029/2019GL083453}
    {}
    {}
  \cvhonor
    {2019}
    {Morse, P. E., A. M. Reading, T. Stål — Well-posed geoscientific visualisation through interactive colour mapping. \textit{Frontiers in Earth Science}, 7, 274. doi:10.3389/feart.2019.00274}
    {}
    {}
\end{cvhonors}
\clearpage
%===============================================================================
% Reports, outreach, code, datasets
%===============================================================================
\cvsection{Reports, outreach, code, datasets}

\begin{cvhonors}
  \cvhonor
    {2025}
    {Reading, A. M., T. Stål, I. Kelly — Ice Sheets and Interactions, Dronning Maud Land and Coats Land [Data set]. \textit{International Federation of Digital Seismograph Networks}. doi:10.7914/j9cq-xd60}
    {}
    {}
  \cvhonor
    {2025}
    {Stål, T., A. M. Reading, S. Kupis, J. Newlands — Casey–Wilkins Adaptable Array [Data set]. \textit{International Federation of Digital Seismograph Networks}. doi:10.7914/SN/Z9\_2021}
    {}
    {}
  \cvhonor
    {2024}
    {Stål, T., F. S. McCormack, A. M. Reading, N. Askey-Doran, J. A. Halpin, M. Lösing — Geothermal heat shapes the Antarctic ice sheet from below. \textit{Frontiers for Young Minds}. doi:10.3389/frym.2023.1178537}
    {}
    {}
  \cvhonor
    {2023}
    {Fuchs, S., F. Neumann, B. Norden, G. Beardsmore, P. Chiozzi, W. Colgan, A. P. Dominguez, M. R. A. Duque, O. M. Ojeda Espinoza, F. Forster, \dots, T. Stål et al. — The Global Heat Flow Database: Update 2023. \textit{GFZ Data Services}. doi:10.5880/fidgeo.2023.023}
    {}
    {}
  \cvhonor
    {2022}
    {Stål, T. — Fast calculation of Ripley's K and L functions on a sphere. \textit{Zenodo}. doi:10.5281/zenodo.6865005}
    {}
    {}
\end{cvhonors}

%===============================================================================
% Academic positions
%===============================================================================
\clearpage
\cvsection{Academic positions}

\begin{cventries}
  \cventry
    {Research Associate in Computational Physics (Geophysics)}
    {ACEAS, University of Tasmania}
    {Tasmania, Australia}
    {2022 -- Present}
    {}

  \cventry
    {Research Fellow in Antarctic Seismology}
    {University of Tasmania}
    {Tasmania, Australia}
    {2021 -- 2022}
    {}

  \cventry
    {Research Assistant}
    {University of Tasmania}
    {Tasmania, Australia}
    {2020 -- 2021}
    {}

  \cventry
    {Casual tutor, demonstrator, and research assistant}
    {University of Tasmania}
    {Tasmania, Australia}
    {2016 -- 2021}
    {}
\end{cventries}

%===============================================================================
% Education
%===============================================================================
\cvsection{Education}

\begin{cventries}
  \cventry
    {PhD, Geophysics, Geology, and Computational Methods}
    {University of Tasmania}
    {Australia}
    {2021}
    {
      \begin{cvitems}
        \item Thesis: \textit{The Antarctic lithosphere revealed by multivariate analysis}.
        \item Supervisors: A. M. Reading, J. A. Halpin, J. M. Whittaker.
        \item Coursework in statistics and computational geophysics.
      \end{cvitems}
    }

  \cventry
    {MSc, Geophysics}
    {University of Copenhagen}
    {Denmark}
    {2015}
    {
      \begin{cvitems}
        \item Thesis: \textit{High-resolution S- and P-wave seismic mapping of near-surface chalk deposits}.
        \item Supervisor: L. Nielsen.
        \item Coursework in geophysics, stochastic modelling, and sedimentary basins.
      \end{cvitems}
    }

  \cventry
    {BSc, Geology}
    {University of Copenhagen}
    {Denmark}
    {2011}
    {
      \begin{cvitems}
        \item Including studies at UNIS, Svalbard.
        \item Thesis: \textit{Processes and development of the Faroe North Slide Complex}.
        \item Supervisors: G. K. Pedersen, T. Nielsen.
        \item Coursework in geology, palaeontology, geomorphology, Quaternary geology, and Arctic marine geology.
      \end{cvitems}
    }

  \cventry
    {AD, Sound Engineering}
    {Piteå Music Conservatory (Luleå University of Technology)}
    {Piteå, Sweden}
    {1998}
    {
      \begin{cvitems}
        \item Coursework in electronics, acoustics, radio broadcasting, music theory, and recording technology.
        \item Including exchange studies at VGIK, Moscow, Russia.
      \end{cvitems}
    }

  \cventry
    {(162.5 ECTS)}
    {Unincorporated studies}
    {}
    {Ongoing}
    {
      \begin{cvitems}
        \item Linguistics (linguistic typology, sociolinguistics, African languages, Nordic studies, pragmatics, grammar).
        \item Geology; climate change; non-Nordic geology (the American Cordillera); Arctic petroleum provinces; geohazards; GIS and remote sensing.
      \end{cvitems}
    }
\end{cventries}

%===============================================================================
% Selected conferences, workshops, hackathons
%===============================================================================
\clearpage
\cvsection{Selected conferences, workshops, and hackathons}

\begin{cvhonors}
  \cvhonor
    {2026}
    {EGU General Assembly, Vienna, Austria — Aq2 and Kq2.}
    {}
    {}
  \cvhonor
    {2025}
    {geoBRIDGE Working Group Hackathon, Perth, Australia — participant and project contributor.}
    {}
    {}
  \cvhonor
    {2025}
    {EGU General Assembly, Vienna, Austria — Seismicity of Denman Glacier: constraints on geometry and dynamics.}
    {}
    {}
  \cvhonor
    {2024}
    {SCAR Open Science Conference, Pucón, Chile — Aq2: refined geothermal heat-flow map of Antarctica; Exploring Atlas Antarktiki; convened three sessions and co-organised collaborations at Bunger Hills.}
    {}
    {}
  \cvhonor
    {2023}
    {EGU General Assembly, Vienna, Austria — Using information entropy to optimise and communicate certainty of continental-scale tectonic models.}
    {}
    {}
  \cvhonor
    {2023}
    {Maydena Subglacial Hydrology and Geology Workshop, Maydena, Australia — subglacial heat budget (conceptual investigation); organised and co-chaired.}
    {}
    {}
  \cvhonor
    {2021}
    {AGU Fall Meeting, New Orleans and online — modelling Antarctic subglacial geothermal heat transfer in higher resolution; co-convened session in Antarctic geothermal heat.}
    {}
    {}
  \cvhonor
    {2021}
    {Australian Earth Sciences Convention (AESC) “Core to Cosmos”, online — data-driven tectonic regionalisation of Antarctica; convened two sessions on data science and machine learning.}
    {}
    {}
  \cvhonor
    {2020}
    {SCAR Open Science Conference 2020: Antarctic – Global Connections, online — towards a reconciled heat-flow map for Antarctica (Aq1); Wilkes Land sector and Aurora Basin geology.}
    {}
    {}
  \cvhonor
    {2020}
    {AGU Fall Meeting, online — multivariate geothermal heat-flow models of Antarctica, towards Aq2.}
    {}
    {}
  \cvhonor
    {2019}
    {Geological Society of Australia SGTSG Meeting, Port Lincoln, Australia — probable geology of the Aurora Basin, East Antarctica.}
    {}
    {}
  \cvhonor
    {2019}
    {XIII International Symposium on Antarctic Earth Sciences (ISAES), Incheon, Korea — linking Antarctic geological observations and geophysical data in a probabilistic space.}
    {}
    {}
  \cvhonor
    {2018}
    {POLAR2018 Open Science Conference, Davos, Switzerland — multi-domain lithospheric model of East Antarctica.}
    {}
    {}
  \cvhonor
    {2018}
    {SCAR TACTical heat-flow workshop, Hobart, Australia — talk and poster on the tectonic structure of Antarctica.}
    {}
    {}
  \cvhonor
    {2018}
    {IEEE Antarctic and Southern Ocean Forum (ASOF), Hobart, Australia — framework for processing multivariate Antarctic Solid Earth data.}
    {}
    {}
  \cvhonor
    {2018}
    {DataTas seminar, Hobart, Australia — Software construction tools make everything so easy!.}
    {}
    {}
  \cvhonor
    {2018}
    {Agile* Subsurface Hackathon, Copenhagen, Denmark — winner of People's Choice Award.}
    {}
    {}
  \cvhonor
    {2017}
    {AGU Fall Meeting, New Orleans, USA — combined constraints on structure and physical properties of East Antarctic lithosphere.}
    {}
    {}
  \cvhonor
    {2015}
    {Workshop on 3D reflection seismic processing using open-source software, Rice University, Houston, Texas.}
    {}
    {}
\end{cvhonors}

%===============================================================================
% Selected non-academic work experience
%===============================================================================
\clearpage
\cvsection{Selected non-academic work experience}

\begin{cventries}
  \cventry
    {GRIT Facility Manager – Instrumentation}
    {University of Tasmania}
    {Tasmania, Australia}
    {2021 -- Present}
    {Instrumentation management and Antarctic fieldwork support.}

  \cventry
    {Geoscientist}
    {Geoneon}
    {Australia}
    {2019 -- 2021}
    {Marine geoscience, GIS, and data compilation.}

  \cventry
    {Field Leader and Geophysicist}
    {COWI}
    {Denmark}
    {2010 -- 2015}
    {MEP, GPR, and TEM fieldwork, and GIS-based survey design and interpretation.}

  \cventry
    {Invited technician and video designer}
    {Al-Harah Theatre}
    {Palestine}
    {2011 -- 2012}
    {}

  \cventry
    {Field Scientist}
    {Norwegian Polar Institute}
    {Svalbard}
    {2010}
    {Supporting biological fieldwork.}

  \cventry
    {Light Designer}
    {Rapid Eye}
    {Denmark}
    {2009 -- 2016}
    {Including the award-winning performance \textit{QUIPROQUO}.}

  \cventry
    {Lighting Designer}
    {STÜ}
    {Estonia}
    {2009 -- 2014}
    {Including touring.}

  \cventry
    {Technical Producer}
    {National Theatre of Ghana}
    {Ghana}
    {2008 -- 2009}
    {}

  \cventry
    {Sound/Video/Light Designer and Technician}
    {The Swedish National Touring Theatre}
    {Sweden}
    {2005 -- 2007 (2011)}
    {}

  \cventry
    {Sound Engineer}
    {Circus Cirkör}
    {Sweden and world tour}
    {2004 -- 2006}
    {}

  \cventry
    {Sound Designer}
    {The Royal Dramatic Theatre}
    {Sweden}
    {2003 -- 2004}
    {}

  \cventry
    {Sound Engineer}
    {Swedish Radio (P1, P3)}
    {Sweden}
    {1997 -- 2004}
    {}

  \cventry
    {Stage Manager}
    {Uppsala Stadsteater}
    {Sweden}
    {1998 -- 2003}
    {}

  \cventry
    {Plumbing Apprentice}
    {Linds Rör AB}
    {Sweden}
    {1993 -- 1996}
    {}
\end{cventries}

%===============================================================================
% Awards & recognition
%===============================================================================
\cvsection{Awards \& recognition}

\begin{cvhonors}
  \cvhonor
    {2023}
    {Royal Society of Tasmania 2023 Doctoral Award.}
    {}
    {}
  \cvhonor
    {2021}
    {Wiley top cited article.}
    {}
    {}
  \cvhonor
    {2020}
    {PhD Award for Outstanding Performance during Postgraduate Studies.}
    {}
    {}
  \cvhonor
    {2019}
    {Progress in Inclusion, Diversity \& Equity.}
    {}
    {}
\end{cvhonors}

%===============================================================================
% Selected memberships and service
%===============================================================================
\cvsection{Selected memberships and service}

\begin{cvhonors}
  \cvhonor
    {2021--}
    {SCAR INSTANT Heat Flow Subcommittee}
    {Co-chair}
    {}
  \cvhonor
    {2023--}
    {International Heat Flow Commission (IHFC), IUGG}
    {Member}
    {}
  \cvhonor
    {2020}
    {Australian Antarctic Research Conference (AARC)}
    {Organising committee}
    {}
  \cvhonor
    {2019}
    {Geological Society of Australia SGTSG Meeting}
    {Organising committee}
    {}
  \cvhonor
    {2017--2020}
    {School of Natural Sciences Inclusion, Diversity \& Equity Committee}
    {Member}
    {}
\end{cvhonors}

%===============================================================================
% Art collaborations (as a scientist)
%===============================================================================
\cvsection{Art collaborations (as a scientist)}

\begin{cvhonors}
  \cvhonor
    {2025}
    {T. Stål, A. Reading, K. Verell, DJ Jon Smeathers — Beaker Street: Seismic Dance Party (Hobart), a collaborative science–art performance translating Denman Glacier seismic data into an immersive audio–visual experience.}
    {}
    {}
  \cvhonor
    {2025}
    {L. Bleach — \textit{An Infrasonorous Archipelago}, Walk\&Talk Biennial “Gestures of Abundance” (Azores), using seismo–volcanic infrasound to create tactile and sonic works.}
    {}
    {}
  \cvhonor
    {2025}
    {M. Nieboer — \textit{Through Our Eyes}, Australian Embassy, Paris; contributed cartographic design and geospatial visualisation.}
    {}
    {}
  \cvhonor
    {2023}
    {L. Bleach — \textit{soft rock monitoring (the slow seismogenic zone)}, EARTHEN. Better Nature, Cement Fondu (Sydney); transmitting vibrational signals from the SYDS urban seismic monitoring site into gallery architecture.}
    {}
    {}
  \cvhonor
    {2021}
    {L. Bleach — \textit{Attenuated ground (the slow seismogenic zone)}, TarraWarra Biennial 2021: Slow Moving Waters; integrating double bass, cast toffee, seismometer, tactile transducers, and slow–earthquake signals.}
    {}
    {}
\end{cvhonors}

%===============================================================================
% Other skills and experience
%===============================================================================
\cvsection{Other skills and experience}

\cvsubsection{Technical skills}
\begin{cvparagraph}
Python (Scikit-learn, OpenCV, TensorFlow, dask, xarray, geopandas), GIS (QGIS, ArcGIS, GDAL/OGR), TeX/LaTeX (TikZ), Bash, Linux, macOS, git, SCons, Madagascar, Seismic Unix, SQL, MAX. Some HTML, C, JavaScript (learning). Experience with Microsoft Office, Adobe Creative Suite, Blender (BlenderGIS), GIMP, Inkscape, Fusion CAD. Working but outdated knowledge of MATLAB, SAC, Emacs Lisp, BASIC.
\end{cvparagraph}

\cvsubsection{Online courses}
\begin{cvparagraph}
Drawing Nature, Science and Culture: Natural History Illustration NHI101x (University of Newcastle, 2020); 3D Earth spring school (University of Kiel, 2021); Presenting Data and Information by Edward Tufte (2025).
\end{cvparagraph}

\cvsubsection{Languages}
\begin{cvparagraph}
English (C2), Swedish (native), Danish (fluent), French (intermediate), Russian (functional), Asante Twi (functional), Norwegian (reading and understanding), Esperanto (inactive), Spanish (learning), Yiddish (learning), Estonian (learning).
\end{cvparagraph}

\cvsubsection{Certificates}
\begin{cvparagraph}
Coastal Skipper Certificate (SWE); Marine HF/VHF Radio Certificate (SWE); firearms licence and hunting certificate (SWE); PADI Advanced Open Water Diver (including dry suit and Nitrox); Wilderness First Aid (AUS); Australian Approved Arrangements Accreditation Biosecurity (AUS); trip leader and terrain-vehicle driver for the Australian Antarctic Division (e.g. Casey 2023).
\end{cvparagraph}

\cvsubsection{Practical skills}
\begin{cvparagraph}
Welding, woodwork, diesel engines, sewing, cooking and baking.
\end{cvparagraph}

\cvsubsection{Other interests}
\begin{cvparagraph}
Music, literature and theatre; particular interest in old boats, especially tugs.
\end{cvparagraph}

\end{document}